\documentclass[../../../include/open-logic-section]{subfiles}

\begin{document}

\olfileid{sth}{replacement}{absinf}
\olsection{প্রতিস্থাপন ও ``পরম অসীমতা''}

এবার \limofsize{} পেছনে রেখে প্রতিস্থাপনের
ন্যায্যতা প্রতিষ্ঠার অন্য এক ধরনের
(অন্তর্নিহিত) চেষ্টা দেখব। এগুলি সেটকে
পর্যায়ে গঠিত হওয়া বস্তু হিসেবে দেখার
ধারণাকে গুরুত্ব দেয়।

পর্যায়ক্রমিক প্রক্রিয়ার প্রথম রূপরেখায়
প্রত্যেক পর্যায়ে কী ঘটে, তা ব্যাখ্যা করতে
কয়েকটি নীতি দিয়েছিলাম: \stageshier,
\stagesord{} এবং \stagesacc। পরে পর্যায়ের
সংখ্যা সম্পর্কে কিছু নীতি যোগ করেছি:
\stagessucc{} বলে সেট গঠনের প্রক্রিয়া
কখনও শেষ হয় না; আর \stagesinf{} বলে
প্রক্রিয়াটি অসীমতম একটি পর্যায়ে পৌঁছায়।

প্রস্তাব করা যুক্তিসংগত যে শেষের দুটি নীতি
আরও বিস্তৃত কোনো নীতি থেকে আসে, যেমন:
\begin{enumerate}
	\item[] \stagesinex. পর্যায় আছে পরমভাবে
	অসীম সংখ্যক; স্তরক্রম যতটা উঁচু হওয়া
	সম্ভব, ততটাই উঁচু।
\end{enumerate}
এটি স্পষ্টতই অনানুষ্ঠানিক নীতি। কিন্তু
সেটের সঞ্চয়ী-পর্যায়ক্রমিক ধারণা থেকে
সরাসরি \emph{অনুসৃত} না হলেও এর সঙ্গে
নীতিটি \emph{সঙ্গতিপূর্ণ} বলেই মনে হয়।
অন্তত \limofsize-র বিপরীতে এটি সেট
পর্যায়ে পর্যায়ে গঠিত হয়—এই ধারণা
বজায় রাখে।

এখন আশা হলো, \stagesinex{} দিয়ে
প্রতিস্থাপনের ন্যায্যতা প্রতিষ্ঠা করা।
কীভাবে হতে পারে দেখি।

\olref[ordinals][idea]{sec}-এ দেখেছি,
এমন একটি সুক্রমবিন্যাস গঠন সহজ,
যা (মূল ভাবনায়) $\omega+\omega$-র
সমরূপ হওয়া উচিত। অন্যভাবে বললে,
পর্যায়ে পর্যায়ে চলা এমন একটি প্রক্রিয়া
সহজেই কল্পনা করতে পারি, যার ক্রমপ্রকার
(মূল ভাবনায়) $\omega+\omega$। তাই
\stagesinex{} মেনে নিলে স্তরক্রমে অন্তত
$\omega+\omega$-তম পর্যায়, অর্থাৎ
$V_{\omega+\omega}$, আছে বলেও মানা উচিত:
স্তরক্রম নিশ্চয়ই \emph{এতদূর} চলতে পারে।

চিন্তাটিকে সাধারণীকরণ করি: যেকোনো
সুক্রমবিন্যাসের জন্য পর্যায়ক্রমিক স্তরক্রম
গঠনের প্রক্রিয়া অন্তত সেই সুক্রমবিন্যাস
পর্যন্ত চলবে। \olref[ordinals][ordtype]{thmOrdinalRepresentation}-কে
\emph{স্বতঃসিদ্ধ} হিসেবে নিলেই এটি
নিশ্চিত করা যায়। তখন যেকোনো
সুক্রমবিন্যাস কোনো ফন নয়মান ক্রমসংখ্যার
সমরূপ হবে। প্রত্যেক ফন নয়মান ক্রমসংখ্যার
স্তরাঙ্ক নিজেই; তাই
\olref[ordinals][ordtype]{thmOrdinalRepresentation}
বলবে, সেটতত্ত্বে কোনো সুক্রমবিন্যাস
বর্ণনা করতে পারলেই সেট গঠনের প্রক্রিয়া
সেই বিন্যাসকে ছাড়িয়ে যেতে হবে।

এই ধারণা \stagesinex-র ফল বলেই
মনে হয়। দুর্ভাগ্যবশত, এটি থেকে
প্রতিস্থাপন পেতে চাইলে একটি সরল
কারিগরি বাধা আসে: প্রতিস্থাপন
\olref[ordinals][ordtype]{thmOrdinalRepresentation}-এর
চেয়ে কঠোরভাবে শক্তিশালী।
(\citet[\S13.2]{Potter2004} এই পর্যবেক্ষণ
করেছেন; \olref[sth][replacement][finiteaxiomatizability]{sec}-এ
আমরা প্রমাণ করব।)

ফলে \stagesinex{}-কে এমনভাবে বুঝতে
হলে যাতে প্রতিস্থাপন পাওয়া যায়,
নীতিটি \emph{শুধু} এই কথা বললে চলবে না
যে স্তরক্রম যেকোনো সুক্রমবিন্যাসকে
ছাড়িয়ে যায়। আরও শক্তিশালী দাবি চাই।
সেই উদ্দেশ্যে \cite{Shoenfield:AST}
ধারণাটির স্বাভাবিক একটি শক্তিশালী রূপ
প্রস্তাব করেছেন: স্তরক্রম কোনো সেটের
সঙ্গে \emph{সহশেষী} নয়।\footnote{গ্যোডেলও
অনুরূপ কথা প্রস্তাব করেছিলেন বলে মনে হয়;
দেখো \citet[p.~223]{Potter2004}। গ্যোডেল
ও \citeauthor{Shoenfield:AST} সম্পর্কে
আলোচনার জন্য দেখো
\citet[90--5]{Incurvati2020}।}
আরও স্পষ্ট করে বললে, $\tau$ যদি
সেটগুলিকে স্তরক্রমের পর্যায়ে পাঠানো
চিত্রণ হয়, তবে কোনো সেট $A$-র
$\tau$-প্রতিচ্ছবি স্তরক্রমকে নিঃশেষ
করে না। অন্যভাবে (ছক আকারে):
\begin{enumerate}
	\item[] \stagescofin. $A$ একটি সেট হলে
	এবং প্রত্যেক $x \in A$-র জন্য
	$\tau(x)$ একটি পর্যায় হলে, এমন একটি
	পর্যায় আছে যা প্রত্যেক $x \in A$-র
	$\tau(x)$-এর পরে আসে।
\end{enumerate}
$\ZF$ থেকে \stagescofin-র উপযুক্ত
আনুষ্ঠানিক রূপ যে অনুসৃত হয়, তা স্পষ্ট।
বিপরীতে, অনানুষ্ঠানিকভাবে যুক্তি দেওয়া যায়
যে \stagescofin{} প্রতিস্থাপনকে ন্যায্য
প্রতিপন্ন করে।\footnote{\stagescofin-র
কোনো আনুষ্ঠানিক রূপ দিয়ে প্রতিস্থাপন
প্রমাণ করা কঠিনতর; কারণ শুধু $\Z$ দিয়ে
পর্যায়গুলি সংজ্ঞায়িত করা যায় না,
তাই \stagescofin{} কীভাবে আনুষ্ঠানিক
করব, তা স্পষ্ট নয়। একটি পথ হলো
\olref[sth][replacement][extrinsic]{sec}-এ
আলোচিত $\LT$-র কোনো বিস্তারে কাজ করা।}
ধরি, $(\forall x \in A)\lexists![y][\phi(x,y)]$।
তাহলে প্রত্যেক $x \in A$-র জন্য
$\sigma(x)$ হলো সেই $y$ যাতে
$\phi(x,y)$, আর $\tau(x)$ হলো প্রথম
যে পর্যায়ে $\sigma(x)$ গঠিত হয়।
\stagescofin{} অনুযায়ী এমন একটি
পর্যায় $V$ আছে যাতে
$(\forall x \in A)\tau(x)\in V$।
এখন প্রত্যেক $\tau(x) \in V$ এবং
$\sigma(x) \subseteq \tau(x)$ বলে
পৃথকীকরণ দিয়ে পাই
$\Setabs{y
\in V}{(\exists x \in A)\sigma(x) = y} = \Setabs{y}{(\exists x \in
A)\phi(x,y)}$।

\begin{prob}
	$\ZF$-এর মধ্যে \stagescofin{}
	আনুষ্ঠানিক রূপ দাও।
\end{prob}

অতএব \stagescofin{} প্রতিস্থাপনকে
সমর্থন করে। আর \stagesinex{} থেকে
\stagescofin{} অনুসৃত হয়, অন্তত
এমন ভাবা সম্ভব। কারণ ধরি,
\stagescofin{} ব্যর্থ। তাহলে স্তরক্রম
কোনো সেট~$A$-র সঙ্গে সহশেষী: এমন
চিত্রণ $\tau$ আছে যে যেকোনো
পর্যায়~$S$-এর জন্য কোনো $x \in A$
পাই যাতে $S \in \tau(x)$। সে ক্ষেত্রে
স্তরক্রমের কথিত ``পরম অসীমতা''কে
ধরার একটি উপায় আছে: $A$-তে
$\tau$ প্রয়োগের বিস্তৃতি স্তরক্রমকে
\emph{নিঃশেষ} করে। তাতে স্তরক্রম
``পরমভাবে অসীম''—এই চিন্তা ক্ষুণ্ণ হয়।
বিপরীতকথনে: \stagesinex{} থেকে
\stagescofin{} অনুসৃত, যা প্রতিস্থাপনের
ন্যায্যতা দেয়।

প্রতিস্থাপনের অন্তর্নিহিত ন্যায্যতা
দেওয়ার এই চেষ্টা সত্যিই আশাব্যঞ্জক।
কিন্তু শেষ পর্যন্ত সফল হয় কি না,
সে সিদ্ধান্ত তোমার উপরই ছাড়ছি।

\end{document}
